\documentclass[../candidatura.tex]{subfiles}
\graphicspath{{\subfix{../imagenes/}}}
\begin{document}

\section{Hipótesis y pregunta de investigación}

\subsection{Hipótesis}

La popusintesíntesis, una práctica artística que consiste en sintetizar sintetizadores de forma popular, es una insistencia en la luthería de instrumentos electrónicos y computacionales para emitir ruidos. Su realización de forma popular, entendida como pública, abierta, reparable, legible, artesanal y en alternativa a las industrias, permite la emergencia de nuevas estrategias de diseño y fabricación de dispositivos, aparatos, máquinas, infraestructuras y metáforas que habilitan la popularización de la producción de ruidos y conocimientos técnicos.

\subsection{Pregunta de investigación}

¿Cómo la popusintesíntesis, entendida como una práctica artística que consiste en sintetizar sintetizadores de forma pública, abierta, reparable, legible y artesanal, genera nuevas estrategias de diseño y de fabricación alternativas a las industrias, orientadas a popularizarla producción de ruidos y conocimientos técnicos en polos tradicionalmente considerados de baja productividad electrónica y computacional?

% en borradores iniciales, la extensión sugerida era de media página

\end{document}